\chapter{Objetivos}

\section{Objetivo General}

Diseñar un agente basado en modelos de lenguaje de gran escala
para el análisis semántico de incidentes de seguridad a partir
de logs heterogéneos, orientado a organizaciones sin equipos
de seguridad especializados.

\section{Objetivos Específicos}

Los objetivos específicos se organizan en tres dimensiones metodológicas: 
\textbf{Estudio}, \textbf{Desarrollo} y \textbf{Validación}, las cuales 
cubren el ciclo completo del proyecto desde la revisión del conocimiento 
existente hasta la evaluación empírica de la propuesta.

% ── ESTUDIO ────────────────────────────────────────────────────
\begin{table}[H]
\centering
\renewcommand{\arraystretch}{1.6}
\begin{tabularx}{\textwidth}{| l | X |}
  \hline
  \multicolumn{2}{|l|}{\textbf{Dimensión: Estudio}} \\
  \hline
  \textbf{OE1} &
  Realizar una revisión sistemática de literatura siguiendo la metodología PRISMA para identificar, analizar y sintetizar las principales contribuciones científicas sobre análisis automatizado de logs de seguridad, detección de anomalías semántica mediante LLMs y arquitecturas de agentes con razonamiento causal, identificando las limitaciones de los enfoques existentes y las brechas que justifican la propuesta. \\
  \hline
  \textbf{OE2} &
  Definir un marco formal que caracterice las propiedades estructurales y semánticas de los logs de seguridad relevantes para la reconstrucción de incidentes, identificando los criterios que debe satisfacer el sistema para mapear eventos a tácticas.
  \\
  \hline
  
\end{tabularx}
\caption{Objetivos específicos correspondientes a la dimensión de Estudio.}
\label{tab:obj_estudio}
\end{table}

\vspace{0.4cm}

% ── DESARROLLO ─────────────────────────────────────────────────
\begin{table}[H]
\centering
\renewcommand{\arraystretch}{1.6}
\begin{tabularx}{\textwidth}{| l | X |}
  \hline
  \multicolumn{2}{|l|}{\textbf{Dimensión: Desarrollo}} \\
  \hline
  \textbf{OE3} &
  Diseñar la arquitectura del agente LLM para el análisis semántico 
  de incidentes de seguridad, especificando sus módulos funcionales, junto con las interacciones, flujos de datos y tecnologías asociadas a cada componente.\\
  \hline
  \textbf{OE4} &
  Implementar un prototipo funcional de un agente basado en modelos de lenguaje que procese logs de seguridad heterogéneos, ejecute mecanismos de razonamiento para la reconstrucción de incidentes y genere reportes interpretables en lenguaje natural.\\
  \hline
\end{tabularx}
\caption{Objetivos específicos correspondientes a la dimensión de Desarrollo.}
\label{tab:obj_desarrollo}
\end{table}

\vspace{0.4cm}

% ── VALIDACIÓN ─────────────────────────────────────────────────
\begin{table}[H]
\centering
\renewcommand{\arraystretch}{1.6}
\begin{tabularx}{\textwidth}{| l | X |}
  \hline
  \multicolumn{2}{|l|}{\textbf{Dimensión: Validación}} \\
  \hline
  \textbf{OE5} &
  Evaluar el desempeño del prototipo mediante experimentos controlados sobre datasets con ground truth, considerando métricas de precisión en la identificación de incidentes, coherencia de las narrativas generadas y comparación con enfoques tradicionales basados en reglas.\\
  \hline
\end{tabularx}
\caption{Objetivo específico correspondiente a la dimensión de Validación.}
\label{tab:obj_validacion}
\end{table}